\documentclass[11pt]{article}
\usepackage[margin=1in]{geometry}
\usepackage{graphicx}
\usepackage{booktabs}
\usepackage{pgfplots}
\pgfplotsset{compat=1.18}
\usepackage{url}
\usepackage[colorlinks=true,citecolor=blue,linkcolor=blue,urlcolor=blue]{hyperref}

% Every number below comes from ../Code/results/ or ../Code/README.md; the
% figure is drawn at build time from ../Code/results/training_metrics.csv.
% Build from paper/: latexmk -pdf -interaction=nonstopmode main.tex

% Keep only rows of one run and phase; rows without a validation measurement
% have an empty val_bpb and are discarded as well. \getthisrow reads the raw
% string, which \thisrow does not inside filter code.
\pgfplotsset{
  select run/.style 2 args={
    unbounded coords=discard,
    x filter/.code={%
      \getthisrow{run}\selrun \getthisrow{phase}\selphase
      \edef\wantrun{#1}\edef\wantphase{#2}%
      \ifx\selrun\wantrun\ifx\selphase\wantphase\else\def\pgfmathresult{nan}\fi
      \else\def\pgfmathresult{nan}\fi
    },
  },
}

\title{nanochat on a Laptop CPU: a Three-Minute End-to-End Example\\for the AI-Human Research OS}
\author{}
\date{September 2026}

\begin{document}
\maketitle

\begin{abstract}
We shrink OpenResearch's default nanochat demo until every stage, from
tokenizer training through pretraining, evaluation, supervised fine-tuning
(SFT), and chat, runs on a laptop-class CPU. With a one-layer,
0.84M-parameter model, a 4{,}096-token vocabulary, and 300 pretraining steps,
the full pipeline, including two probe runs, finishes in 2~min~57~s on an Intel
i7-12700, using at most 1.0~GB of RAM per process. Validation bits per byte
(bpb) falls from 3.456 to 2.288 during pretraining and from 2.289 to 2.146
during SFT. Doubling the Muon matrix learning rate lowers the final bpb by
0.012, and halving the vocabulary raises it by 0.096. The model produces no
useful answers: the contribution is a runnable, deterministic first example
for the Research OS, not a research result.
\end{abstract}

\section{Introduction}

The AI-Human Research OS organizes research as plain files that a human and
code agents share. Its users need one example project they can run end to end
on their own machine, to see the idea, code, figures, results, and paper live
together. OpenResearch~\cite{alphaxiv2026openresearch} ships such a demo: it
trains nanochat~\cite{karpathy2025nanochat}, a compact pipeline for building a
chat model, at depth~6 with 73.5M parameters for 5{,}000 pretraining and
1{,}500 SFT steps on an Apple M3 Max. Its pretraining alone would take over
nine hours on the CPU used here.

This project ports that demo and scales it down until the whole pipeline takes
about three minutes on a CPU. It contributes (i)~a one-command, staged,
CPU-only port whose code changes are documented; (ii)~measurements of the
time, memory, and learning curves of every stage; and (iii)~the demo's two
probes, a larger matrix learning rate and a smaller vocabulary, run at this
scale.

\section{Setup}

\paragraph{Code.} The port starts from OpenResearch commit \texttt{27cb342},
path \texttt{demo/nanochat}, with the demo's overrides applied. The code
changes are small: \texttt{torch.compile} becomes switchable and is off on CPU,
because the Windows CPU backend needs a C++ compiler; SFT data can be
downloaded as an evenly spaced subsample; and CORE prompts are cropped to the
rotary cache length. A staged runner executes each stage in under a minute
and logs its exit code, wall time, and peak memory.

\paragraph{Recipe.} Table~\ref{tab:recipe} compares the demo's recipe with
this example. The data-to-parameter ratio of pretraining stays close to the
demo's (3.95 versus 3.5 tokens per scaling parameter). Pretraining uses one
training shard and the pinned validation shard of
\texttt{climbmix-400b-shuffle}~\cite{karpathy2026climbmix}. SFT uses the
demo's mixture of SmolTalk~\cite{allal2025smollm2}, MMLU~\cite{hendrycks2021mmlu},
and GSM8K~\cite{cobbe2021gsm8k}, subsampled to about 1/128 of every split.
Matrix parameters are trained with Muon~\cite{jordan2024muon}, as in nanochat.
SFT starts at 0.1 of the inherited learning rate: at the demo's 0.8, SFT
validation bpb spiked from 2.29 to 2.63 and ended above its start.

\begin{table}[t]
\centering
\small
\begin{tabular}{lll}
\toprule
Knob & Demo & This example \\
\midrule
Model & depth 6, width 384, 73.5M params & depth 1, width 64, 0.84M params \\
Tokenizer & 32{,}768 vocab, 2B chars & 4{,}096 vocab, 50M chars \\
Context / batch & 512 / 16{,}384 tokens & 256 / 4{,}096 tokens \\
Pretraining data & 8 shards (${\sim}0.8$~GB) & 1 train + 1 val shard (${\sim}184$~MB) \\
Pretraining steps & 5{,}000 (81.9M tokens) & 300 (1.23M tokens) \\
SFT data & full splits (${\sim}980$~MB) & ${\sim}1/128$ of each split (${\sim}14$~MB) \\
SFT steps / LR fraction & 1{,}500 / 0.8 & 40 / 0.1 \\
Base evaluation & CORE + bpb + samples & bpb + samples \\
\bottomrule
\end{tabular}
\caption{Recipe of the OpenResearch demo and of this example.}
\label{tab:recipe}
\end{table}

\paragraph{Hardware and metric.} All runs use Windows~11 on an Intel
i7-12700 (12 cores, 20 threads) with 32~GB of RAM and no GPU, in eager mode.
The metric is validation bits per byte (bpb), measured every 25 pretraining
steps on 65{,}536 validation tokens and every 10 SFT steps on 32{,}768 tokens.
Training is deterministic on this machine: two fresh runs gave identical
validation bpb to six decimals.

\section{Results}

\paragraph{Learning curves.} Pretraining lowers validation bpb from 3.4564 to
2.2882 in 300 steps (Figure~\ref{fig:curves}). The base evaluation on
16{,}384 tokens per split gives 2.2947 bpb on training data and 2.3622 on
validation data. SFT then lowers validation bpb from 2.2892 to 2.1456 in 40
steps. Asked ``What is the capital of France?'', the SFT model replies with
nonsense (\texttt{ATedning the date that is - min B}, whitespace collapsed),
as expected at 0.84M parameters; the demo's 73.5M model answered ``Paris''.

\begin{figure}[t]
\centering
\begin{tikzpicture}
\pgfplotstableread[col sep=comma]{../Code/results/training_metrics.csv}\metrics
\begin{axis}[
  width=0.82\linewidth, height=6.2cm,
  xlabel={Pretraining step}, ylabel={Validation bpb},
  xmin=0, xmax=300, grid=major, grid style={gray!20},
  legend pos=north east, legend cell align=left,
]
\addplot+[thick, mark=*, mark size=1.2pt, select run={baseline}{base}]
  table[x=step, y=val_bpb]{\metrics};
\addlegendentry{Baseline (vocab 4{,}096)}
\addplot+[thick, mark=square*, mark size=1.2pt, select run={probe_lr2x}{base}]
  table[x=step, y=val_bpb]{\metrics};
\addlegendentry{Probe 1: matrix LR $\times 2$}
\addplot+[thick, mark=triangle*, mark size=1.4pt, select run={probe_vocab2048}{base}]
  table[x=step, y=val_bpb]{\metrics};
\addlegendentry{Probe 2: vocab 2{,}048}
\end{axis}
\end{tikzpicture}
\caption{Validation bpb during pretraining for the baseline and both probes,
drawn from \texttt{Code/results/training\_metrics.csv}.}
\label{fig:curves}
\end{figure}

\paragraph{Probes.} Both probes repeat the 300-step schedule with one change
(Table~\ref{tab:probes}). Doubling the Muon matrix learning rate is neutral
through step 100 and then slightly better, ending 0.012 bpb lower. The
2{,}048-token vocabulary is worse at every step; the gap narrows from 0.23 to
0.10 bpb but does not close. The demo's own probe results are not part of the
copied evidence, so whether these directions match the demo's is not checked.

\begin{table}[t]
\centering
\small
\begin{tabular}{rrrrrr}
\toprule
Step & Baseline & Probe 1 (LR $\times 2$) & $\Delta$ & Probe 2 (vocab 2{,}048) & $\Delta$ \\
\midrule
0   & 3.4564 & 3.4564 & 0.0000  & 3.6833 & $+0.2269$ \\
50  & 2.8124 & 2.8155 & $+0.0031$ & 2.9829 & $+0.1704$ \\
100 & 2.5360 & 2.5358 & $-0.0002$ & 2.6756 & $+0.1396$ \\
200 & 2.3436 & 2.3337 & $-0.0100$ & 2.4474 & $+0.1038$ \\
300 & 2.2882 & 2.2761 & $-0.0121$ & 2.3840 & $+0.0958$ \\
\bottomrule
\end{tabular}
\caption{Validation bpb of the probes against the baseline;
$\Delta$ = probe $-$ baseline, lower is better.}
\label{tab:probes}
\end{table}

\paragraph{Cost.} On a first run with an empty data cache, the whole pipeline
takes 2~min~57~s, of which the baseline stages take 1~min~29~s. Pretraining
takes 42.3~s, SFT 20.5~s, and the two probes 47.1~s and 38.4~s. The largest
process peaks at 1.0~GB of RAM (the data download); training stays at or below
0.94~GB. The run downloads about 0.20~GB of data and uses about 0.9~GB of disk
inside the project, mostly for the Python environment.

\section{Limitations}

This is a demonstration of the pipeline, not a model: the chat replies are
nonsense, and the bpb values and probe deltas say nothing about larger models.
Each probe uses one seed; runs are deterministic, so the deltas are not
run-to-run noise, but seed-to-seed variance was not measured, and probe~1's
$-0.012$~bpb may not be meaningful. Probe~2 changes more than the vocabulary:
the smaller model has 0.44M instead of 0.84M parameters and sees fewer bytes
of text per 4{,}096-token step. The automatic weight-decay scaling gives
$\lambda = 7.0$ at this size, far outside the range the scaling rules were
tuned for. SFT sees mostly short conversations, because the loader skips
those longer than the 256-token context. Only this Windows machine was tested.

\section{Reproducibility}

From the project's \texttt{Code/nanochat/} directory, \texttt{bash
runs/runcpu\_small.sh all} runs every stage and writes
\texttt{Code/results/summary.json} and
\texttt{Code/results/training\_metrics.csv}; the timings above come from the
first run recorded in \texttt{Code/README.md}. This paper builds from
\texttt{paper/} with \texttt{latexmk -pdf main.tex} and draws
Figure~\ref{fig:curves} from the tracked CSV at build time.

\bibliographystyle{plain}
\bibliography{references}

\end{document}
